\begin{description}
\item[Est2を作るブロック] MCMCサンプルをデータフレーム化したものの中から，状態$\mu$に関するデータだけを抜き出したものを作る。
\item[Est2missを作るブロック] 欠損値を補間するための推定値を，MCMCサンプルから取り出した欠損値補間だけのデータセットを作る。
\item[plot用の関数] プロットはこの後にも行いますので，もう関数を作ってまとめてしまうことにします。もとのデータセット\verb|fullDays|と，状態の推定値データセット\verb|Est|，欠損値の推定値データセット\verb|MissEst|を引数にとる関数です。
\begin{description}
	\item[関数3行目] まずもとデータに行番号を\verb|ID|として変数化します。
	\item[関数4行目] 状態の推定値は，すべての行について推定しているはずですので，行番号をキーに結合\verb|left_join|します。
	\item[関数5行目] もとの観測値が実測値なのか欠損値なのかを表現する変数\verb|FLG|を用意しておきます。これらはすべて，一時的なオブジェクト\verb|tmp|に納められます。
	\item[関数6行目] 欠損値のインデックスを表す変数\verb|misJ|をつくり，カウントを1に設定しておきます。
	\item[関数7-9行目] 先ほどの一次オブジェクトにいまから体重の推定値を入れていきますが，この推定値は実測値で得られている時は不要で\verb|NA|になるはずですので，いったんすべてに\verb|NA|を代入しています。
	\item[関数10-21行目] オブジェクト\verb|tmp|を一行ずつ見ていって，もし\verb|FLG|変数が欠損値であることを教えてくれたら，\verb|misJ|番目の推定値を代入します。代入が終わったらカウンタを1つ追加しておきます。欠損値でなければ，もとの体重をそのまま移し入れます。ちなみにUやLがついているのは95\%区間の上限と下限なのですが，実測値の場合は推定ではありませんのですべて同じ数字になります。
	\item[関数22行目] 作ったオブジェクト\verb|tmp|を戻り値として返します。
\end{description}
